\subsection{La estructura básica del ahora: \textit{duración} como condición de posibilidad de la temporalidad de las vivencias}\label{sec:estructura-ahora}

\begin{quote}
``El tiempo es rígido y, con todo, el tiempo fluye''
\end{quote}

\cite[§31, p.~84]{husserl2002}

El ahora, como rasgo esencial, posee un carácter de presencia total. El
ahora se presenta como el momento que siempre está marcando el reloj.
Cuando hablamos del ahora, solemos pensar en este como el momento más
``real'', incluso, como lo ``único que existe''. Dejamos fuera de él lo
que ya ha pasado o todavía no pasa, ya que se diferencian claramente
respecto al ahora. Por más que el ahora irreflexivamente se presente
como el punto temporal en el que siempre estamos, podemos afirmar que
todo ahora posee un horizonte temporal en relación necesaria con lo
pasado y futuro. Si no fuera así, viviríamos en un \emph{eterno
presente,} lo cual impediría el despliegue temporal necesario para la
multiplicidad de vivencias conscientes. El paso de un ``ahora'' a otro
sería un salto abismal, como la infinidad de decimales que hay entre el
1 y el 2. Si todo ``ahora'' que ocurre no demanda un horizonte temporal
de pasado y futuro, no sería posible distinguir la multiplicidad de
puntos temporales que se enumeran en el tiempo objetivo. No podemos ver
la hora en el reloj sin esta continuidad entre el momento que acaba de
pasar y el que va a suceder. Para poder ``dar'' la hora en cada momento
marcado numéricamente, ya se debe de suponer una duración continua que
el tiempo objetivo no muestra en el paso de un instante a otro, sino que
presupone.

Las consecuencias de la comprensión discreta del tiempo se pueden
encontrar en los estudios de Franz Brentano, maestro de Husserl, quien,
para explicar cómo podemos tener una noción de la temporalidad de
nuestros actos perceptivos recurre a la fantasía para explicar cómo
podemos sintetizar, unir o \emph{asociar} distintos momentos de
captación \cite[p.~54]{mensch2010}. Mediante el acto de la fantasía podemos
``ampliar nuestra intuición temporal'' y proyectar vivencias al pasado o
futuro de diferentes fenómenos psíquicos:

\begin{quote}
``\ldots Toda clase de deseos y actos de voluntad, así como las
convicciones y opiniones de todo tipo, y una amplia gama de
representaciones que tienen su origen en la fantasía. De todos los
fenómenos psíquicos, únicamente las sensaciones --y ni siquiera todas
ellas-- permanecen dentro del ámbito de lo mensurable'' \cite[p.~52]{brentano2009}\footnote{Traducción de la versión en inglés}.
\end{quote}

Brentano recurre a explicar los procesos asociativos de vivencias a
través del acto de la fantasía para explicar la relación de la
percepción presente con vivencias pasadas y futuras: ``Brentano, por lo
tanto, atribuye la conjunción a un proceso de `asociación primordial'
por el cual la fase pasajera de una apariencia se une a la fase ahora
como una `idea de fantasía'\,'' \cite[p.~188]{morrison1970}. Husserl, como
veremos, no busca recurrir a otros actos conscientes para explicar la
asociación de múltiples vivencias, sino describir una estructura
temporal que atraviesa todo acto. Esto hace que la temporalidad no sea
producto de la percepción consciente, sino que esta última necesite de
una estructura temporal para articularse.

Los exámenes fenomenológicos de la conciencia interna del tiempo
muestran cómo cada vivencia actual se extiende en distintas fases que le
dan una continuidad temporal y entrelazan con el horizonte de vivencias
pasadas y futuras. Dicha continuidad de la temporalidad interna de la
conciencia explica la duración temporal de cada acto y de la vida
consciente en su conjunto sin recurrir a una noción discreta de tiempo
que comprende al ahora como un instante aislado en una serie lineal de
sucesiones. Brough comenta que la captación de la duración del ahora
descrito en los exámenes fenomenológicos del tiempo de la conciencia
demuestra que incluso nuestro lapso o duración de tiempo más ``corta''
posee más fases que el instante puntual en cada momento sustituido por
uno nuevo:

\begin{quote}
Así que incluso en el momento presente ocurre una cierta indecidibilidad
entre uno y muchos, y con la comprensión metafísica del tiempo como una
sucesión lineal de momentos del ahora, el momento realmente presente se
encuentra acompañado por miradas de otros momentos que lo preceden o lo
siguen. Cada uno puede ser una unidad considerada en sí misma, pero el
punto es que ninguno de ellos por sí mismo podría componer el
tiempo---el tiempo exige que haya muchas de tales unidades---y así se
compromete el valor de la mera presencia de un sólo momento. La
comprensión del momento por parte de Husserl es quizás menos
pirotécnica, pero ciertamente es sutil e introduce en el ahora un
sentido más profundo de indecidibilidad \cite[p.~517]{brough1993}.
\end{quote}

En contraste con el tiempo objetivo, sujeto a una discreción lineal de
instantes que pueden ser numéricamente medidos, Husserl expone que la
temporalidad consciente no se reduce a una serie de instantes, sino que
consiste en una duración que implica siempre la distinción de fases que
corresponden a lo ya pasado, a lo presente y lo anticipado. Limitaciones
de una mirada discreta del tiempo se muestran al reflexionar sobre
cualquier vivencia, por ejemplo, el oír una melodía o hablar. Cuando
recuerdo dichas vivencias, no recuerdo la suma de notas o de sílabas
pronunciadas sucesivamente, sino que recuerdo la unidad desplegada de la
melodía o las palabras. Originalmente, el tiempo se da en esta unidad
constitutiva de toda experiencia. La comprensión discreta del tiempo, si
bien puede lograr una exactitud y utilidad para propósitos de medición o
cálculo de fenómenos, es una aproximación artificial, no intuitiva, al
tiempo, ya que no da cuenta del fenómeno original de la continuidad
temporal con la que captamos objetos que, bajo una determinada mirada
científica, podemos medir\footnote{Para Husserl, un rasgo esencial de la
  intuición del tiempo que este se dé en una duración: ``Lo que sí
  pertenece a la esencia de la intuición de tiempo es, en cambio, el ser
  a cada punto de su duración (que nosotros podemos convertir
  reflexivamente en objeto) conciencia de \emph{lo que acaba de ser} y
  no mera conciencia del punto de ahora del objeto que aparece durando''
  (Husserl 2002, §12, 54). La duración, en contraste con la mera
  presencia, nos abre el camino a examinar cómo lo que está presente se
  constituye a partir de lo que no lo está, ya sea porque \emph{ya no
  es} o porque \emph{todavía no lo es.}}.

Se exige la distinción de tres momentos temporales: el ``ahora
puntual'', el ``ahora-reciente'' y el ``ahora-adveniente''. Estos tres
momentos se constituyen en la misma conciencia inmanente, la cual el
fenomenólogo examina en las vivencias y puede distinguir en la
conciencia una duración que no es discreta sino continua. Son momentos
enteramente abstractos y se les introduce en la descripción de la forma
en la que se articula la duración mínima. La duración inmanente se
muestra a la mirada del fenomenólogo como un flujo continuo de
vivencias, como una duración de duraciones \cite[p.~120]{husserl2002}. Así, los
tres momentos que se distinguen son enteramente abstractos: cada ahora
puntual está a la vez abierto hacia adelante y adherido hacia atrás.
Correlativamente a estos fenómenos, el fenomenólogo distingue los
momentos de la propia conciencia de tiempo inmanente. Esta
pre-estructura noético-noemática, de los momentos del ahora puntual, el
ahora adveniente y el ahora-reciente es anterior a toda constitución de
tiempo objetivo. Yo puedo \emph{medir} la cantidad de tiempo (objetivo)
que tarda un lápiz de caer al piso desde la mesa. Sin embargo, esta
medición \emph{supone la vivencia} de la percepción de la caída del
lápiz que se despliega en una duración temporal. Así, el ``ahora''
cumple una función tanto en el tiempo objetivo como en la temporalidad
interna de la conciencia. Yo puedo saber que ahora son las 11:00 am
porque el reloj está sucesivamente contando la serie de instantes
pasados y futuros para poder, en cada momento, apuntar al ahora. Sin
embargo, la conciencia del ahora se distingue porque \emph{no marca} un
ahora puntual, sino que todo momento \emph{vivido} implica una
estructura tripartita de duración. De este modo, la distinción de las
fases de la conciencia del ahora, sin embargo, no debe ser entendida
como una \emph{separación} sucesiva de instantes discretos: se dan más
bien \emph{como una unidad en cada} vivencia:

\begin{quote}
Del fenómeno decursivo sabemos que es una continuidad de constantes
mudanzas que forma una unidad inseparable; inseparable en trechos que
pudiesen existir por sí, e indivisible en fases que pudiesen existir de
por sí, en puntos de la continuidad. Los fragmentos que nosotros
abstractamente destacamos sólo pueden existir en el decurso íntegro, y
lo mismo las fases, los puntos de la continuidad decursiva (Husserl
2002, §10, 51).
\end{quote}

En otras palabras, la conciencia puede reflexionar sobre su propio
discurrir temporal, que muestra su unidad debido a la estructura
temporal encontrada en el ahora y que se distingue del tiempo objetivo
en el que discurren los objetos captados en la actitud natural. Así,
puede cambiar de objeto percibido, de acto mediante el cual lo aprehende
y su contenido y, sin embargo, su unidad es asegurada por esta cadena
estructurada de fases temporales que conectan toda vivencia pasada
presente y futura de la cual se pueda tener experiencia. La distinción
de las fases temporales de la conciencia del ahora muestra que el tiempo
interno de la conciencia es una temporalidad compleja, que no se
constituye a partir de instantes homogéneos en una sucesión lineal, sino
que, toda duración temporal implica una distinción de lo ya pasado, lo
presente y lo adveniente. Esto se corrobora al examinar cómo los
exámenes de la conciencia del tiempo se modifican con el avance de sus
exploraciones fenomenológicas, ya que las funciones de las fases de
retención y protención cobran una mayor relevancia ``Los cambios en el
diagrama\footnote{El diagrama se expondrá en la siguiente parte.} y su
descripción reflejan el cambio de enfoque de Husserl hacia los dos
sistemas de protenciones y retención y lejos de la zona de
actualización'' \cite[p.~135]{rodemeyer2003}.